\chapter{Trade Policy}\label{ch:b3:trade-policy}

A tariff is the oldest instrument of economic policy, and for half a century most
countries cut theirs. This chapter asks what a tariff does when the country that
levies it is large enough to move world prices, why that creates a temptation and
a trap, how quotas and export restraints differ from tariffs, how to measure
the protection an industry really gets, when a subsidy can shift profit to a home
firm, why governments levy tariffs at all, and what the rules of the trading
system do to the temptation. It ends with the evidence: who paid the
tariffs of 2018--2019, what import competition did to the workers exposed to it,
and whether protecting an infant industry ever made it grow up.

\begin{openingchart}{Applied tariff rate, mean over all products weighted by
imports, 1988--2022 (percent, times-ten scale). The United States value for 2019 (13.8~\%)
is a break in the source series and is not plotted.}{World Bank, World
Development Indicators (TM.TAX.MRCH.WM.AR.ZS)}
\begin{tikzpicture}
\begin{semilogyaxis}[omtimesten, width=0.7\linewidth, height=0.45\linewidth, xmin=1987, xmax=2023,
  ymin=0.7, ymax=80, ylabel={Percent}, legend columns=3,
  legend style={at={(0.5,1.03)}, anchor=south, font=\footnotesize}, unbounded coords=jump]
  \addplot[omDef, thick, mark=none] table[x=year, y=USA] {figdata/bachelor-3/11-trade-policy/out/tariffs.csv};
  \addplot[omThm, thick, mark=none] table[x=year, y=EUU] {figdata/bachelor-3/11-trade-policy/out/tariffs.csv};
  \addplot[black, thick, mark=none] table[x=year, y=JPN] {figdata/bachelor-3/11-trade-policy/out/tariffs.csv};
  \addplot[omDef, thick, dashed, mark=none] table[x=year, y=CHN] {figdata/bachelor-3/11-trade-policy/out/tariffs.csv};
  \addplot[omThm, thick, dashed, mark=none] table[x=year, y=IND] {figdata/bachelor-3/11-trade-policy/out/tariffs.csv};
  \addplot[black, thick, dashed, mark=none] table[x=year, y=BRA] {figdata/bachelor-3/11-trade-policy/out/tariffs.csv};
  \legend{United States, European Union, Japan, China, India, Brazil}
\end{semilogyaxis}
\end{tikzpicture}
\end{openingchart}

% figdata: tariffs_points, tariffs_dropped
India's average applied tariff fell from 56.4~\% in 1990 to 4.6~\% in 2022, China's
from 32.2~\% in 1992 to 2.2~\%, Brazil's from 31.9~\% to 7.3~\%. The United States, the
European Union and Japan started below 5~\% and ended below 2~\%. The tariff
increases of 2018--2019, large on the goods they targeted, are hardly visible in an
average over all imports.

\begin{recall}
The grades 10--12 volume analysed a tariff in a small country: its incidence,
its revenue and its deadweight loss, and the Year 1 volume measured the rent
created by a quota. Nash equilibrium, the prisoner's dilemma and repeated games
are the Year 2 volume's. \Cref{ch:b3:credible-empirical-work} graded evidence.
\end{recall}

\begin{notation}[Symbols of this chapter]\label{not:b3:trade-policy:symbols}
$t$ is an ad valorem tariff; $P^*$ is the world price received by foreign exporters
and $P = (1 + t)P^*$ the domestic price. $M(P)$ is the home country's import demand
and $X^*(P^*)$ the foreign export supply; $\varepsilon^* = X^{*\prime}P^*/X^*$ is its elasticity
(here not the trade elasticity of \cref{ch:b3:gravity}).
\end{notation}

\section{A tariff in a large country}\label{sec:b3:trade-policy:large}

A small country takes the world price as given: a tariff raises the domestic
price by its full amount and costs the country its deadweight loss. A large
country's imports are a large part of the world market. When it cuts its
purchases, the world price falls, and part of the tariff is paid by foreign
exporters.

\begin{definition}[Terms-of-trade effect; optimal tariff]\label{def:b3:trade-policy:optimal-tariff}
\index{terms-of-trade effect}\index{optimal tariff}
The \emph{terms-of-trade effect} of a tariff is the gain to the importing country
from the fall in the world price of its imports, $-M\,dP^*$. The \emph{optimal
tariff} is the tariff that maximises the importing country's welfare, given the
foreign export supply and no retaliation.
\end{definition}

\begin{proposition}[The welfare effect of a tariff]\label{prop:b3:trade-policy:welfare}
With competitive markets, home welfare $W$ (consumer and producer surplus plus
tariff revenue) changes with the tariff by
\[
  dW = -M\,dP^* + tP^*\,dM .
\]
The first term, the terms-of-trade gain, is positive; the second, the efficiency
loss, is negative for $t > 0$ and zero at $t = 0$. A small tariff therefore raises
the welfare of a country that faces an upward-sloping foreign export supply.
\end{proposition}

\begin{proof}
Consumer plus producer surplus falls by $M\,dP$ when the domestic price rises by
$dP$, and revenue is $tP^*M$. So $dW = -M\,dP + d(tP^*M) = -M\,dP + P^*M\,dt + tM\,dP^* +
tP^*\,dM$. With $P = (1 + t)P^*$, $dP = P^*dt + (1 + t)dP^*$; substituting, the $P^*M\,dt$
terms cancel and $-M(1 + t)dP^* + tM\,dP^* = -M\,dP^*$. A higher tariff lowers $P^*$, so
$-M\,dP^* > 0$, and lowers $M$, so $tP^*dM \le 0$, with equality at $t = 0$.
\end{proof}

\begin{omfigure}
\begin{tikzpicture}[econ, x=0.12cm, y=0.2cm]
  \fill[omThm!25] (0,17.78) rectangle (31.11,20);
  \fill[omDef!25] (31.11,22.22) -- (40,20) -- (31.11,17.78) -- cycle;
  \draw[econ axes] (0,0) -- (62,0) node[right] {$M$};
  \draw[econ axes] (0,0) -- (0,32) node[above] {$P$};
  \draw[curve, omDef] (0,30) -- (60,15) node[right, font=\scriptsize] {$M(P)$};
  \draw[curve, omThm] (0,10) -- (60,25) node[right, font=\scriptsize] {$X^*(P^*)$};
  \draw[curve, omThm, dashed] (0,12.5) -- (56,30) node[right, font=\scriptsize] {$(1 + t)P^*$};
  \draw[guide] (0,22.22) node[left, font=\scriptsize] {$P$} -- (31.11,22.22) -- (31.11,0) node[below, font=\scriptsize] {$M_t$};
  \draw[guide] (0,20) node[left, font=\scriptsize] {$P_0$} -- (40,20) -- (40,0) node[below, font=\scriptsize] {$M_0$};
  \draw[guide] (0,17.78) node[left, font=\scriptsize] {$P^*$} -- (31.11,17.78);
\end{tikzpicture}
\par\smallskip{\small A tariff in a large country, in the import market. The tariff
drives a wedge between the domestic price $P$ and the world price $P^*$. The
rectangle is the terms-of-trade gain, paid by foreign exporters; the triangle is
the efficiency loss.}
\end{omfigure}

\begin{proposition}[The optimal tariff]\label{prop:b3:trade-policy:optimal}
If the \omterm{def:b3:trade-policy:optimal-tariff}{optimal tariff} is interior, it equals the inverse of the elasticity of
foreign export supply at the optimum: $t^{\mathrm{opt}} = 1/\varepsilon^*$.
\end{proposition}

\begin{proof}
At the optimum $dW = 0$: $-M\,dP^* + tP^*\,dM = 0$. Market clearing, $M = X^*(P^*)$, gives
$dM = X^{*\prime}dP^*$. Hence $-M + tP^*X^{*\prime} = 0$, so $t = M/(P^*X^{*\prime}) = X^*/(P^*X^{*\prime}) =
1/\varepsilon^*$.
\end{proof}

A small country faces an infinitely elastic export supply: its \omterm{def:b3:trade-policy:optimal-tariff}{optimal tariff} is
zero. A large one gains from a tariff only at the expense of its partners: the
terms-of-trade gain is their loss, and the world as a whole loses the two
efficiency triangles.

\section{Retaliation: the tariff game}\label{sec:b3:trade-policy:game}

If each country gains from its own tariff and loses from the other's, the choice
of tariffs is a game. With payoffs measured from free trade, a typical case is
the one below.

\begin{omfigure}
\begin{tikzpicture}[econ, font=\footnotesize]
  \node[payoff] at (0.8,0.4) {0, 0};
  \node[payoff] at (2.4,0.4) {$-6$, 4};
  \node[payoff] at (0.8,-0.4) {4, $-6$};
  \node[payoff] at (2.4,-0.4) {$-2$, $-2$};
  \node[anchor=east] at (-0.05,0.4) {free trade};
  \node[anchor=east] at (-0.05,-0.4) {tariff};
  \node at (0.8,1.1) {free trade}; \node at (2.4,1.1) {tariff};
  \node[font=\small\bfseries] at (1.6,1.7) {Foreign};
  \node[font=\small\bfseries, rotate=90] at (-2.0,0) {Home};
\end{tikzpicture}
\par\smallskip{\small The tariff game. Each cell gives Home's payoff, then
Foreign's. Whatever the other does, a tariff pays: the equilibrium is mutual
tariffs, worse for both than free trade.}
\end{omfigure}

It is a prisoner's dilemma. In a single round, both impose tariffs; in a repeated
game, the threat of retaliation can sustain free trade if the countries value
the future enough (\cref{exo:b3:trade-policy:7}). Trade agreements can be read
this way: commitments that make retaliation predictable, and the gain from
deviating short-lived.

\section{Quotas and export restraints}\label{sec:b3:trade-policy:quotas}

\begin{definition}[Import quota; voluntary export restraint]\label{def:b3:trade-policy:quota}
\index{import quota}\index{voluntary export restraint}
An \emph{import quota} limits the quantity of a good that may be imported. A
\emph{voluntary export restraint} (VER) is a quota administered by the exporting
country, which allocates the right to export among its own firms.
\end{definition}

\begin{proposition}[Tariffs and quotas]\label{prop:b3:trade-policy:quota}
Under perfect competition, a quota equal to the imports under a tariff produces
the same domestic price, quantities and deadweight loss; only the owner of the
gap between domestic and world price --- the government's revenue under the
tariff, the quota rent under the quota --- differs. If the domestic industry is a
monopoly, the two are not equivalent: the quota gives the monopolist market power
that the tariff does not.
\end{proposition}

\begin{proof}
Under competition, with imports fixed at $\bar M = M(P_t)$, the domestic market
clears where $D(P) - S(P) = \bar M$, which is the price $P_t$ of the tariff, so every
quantity is the same. With a domestic monopolist and a tariff, imports are
available at $P_t$ in any amount: the monopolist cannot sell above $P_t$ and, as a
price taker there, produces where marginal cost equals $P_t$. With a quota $\bar M$,
the monopolist faces the residual demand $D(P) - \bar M$, sets its own marginal revenue
equal to marginal cost, and chooses a higher price and lower output than under the
tariff (\cref{exo:b3:trade-policy:11}).
\end{proof}

With a VER, the quota rent goes to the foreign exporters, who sell less at a higher
price: the importing country loses the rent as well as the triangles. That is why
exporting firms have accepted such restraints.

\section{Effective protection}\label{sec:b3:trade-policy:erp}

\begin{definition}[Effective rate of protection]\label{def:b3:trade-policy:erp}
\index{effective rate of protection}
The \emph{effective rate of protection} (ERP) of an industry is the proportional
increase in its value added per unit, caused by the tariffs on its output and on
its inputs, relative to value added at world prices.
\end{definition}

Let a good sell at the world price $P$ and use inputs worth $aP$ at world prices.
Value added is $P(1 - a)$. With a tariff $t_f$ on the good and $t_i$ on the inputs, it
becomes $P[(1 + t_f) - a(1 + t_i)]$, so
\begin{equation}\label{eq:b3:tp:erp}
  \mathrm{ERP} = \frac{t_f - a\,t_i}{1 - a} .
\end{equation}

\begin{measure}[Effective protection]
The ERP exceeds the tariff on the good when inputs are taxed less, and becomes
negative when inputs are taxed and the output is not. A tariff on steel is then a
tax on carmakers and every other steel user: they pay more for an input and
compete with foreign rivals who do not. Tariff schedules that rise with the stage
of processing --- low on raw materials, high on finished goods --- give the final
stage an effective protection several times its nominal tariff
(\cref{exo:b3:trade-policy:3}).
\end{measure}

\section{Strategic trade policy}\label{sec:b3:trade-policy:strategic}

The \omterm{def:b3:trade-policy:optimal-tariff}{optimal tariff} works through the world price. A second argument for
intervention works through the profits of firms with market power. When a few
firms share a world market, each earns a rent, and a government that can shift
that rent towards its own firm may gain even though the world loses.

\begin{definition}[Strategic trade policy]\label{def:b3:trade-policy:strategic}
\index{strategic trade policy}
\emph{Strategic trade policy} is a subsidy or a tax on a country's exports in a
market served by a few firms, chosen to change the rivals' behaviour so that
profit moves to the home firm.
\end{definition}

Take a Home firm and a Foreign firm that sell the same good only in a third
market, with inverse demand $P = a - q_h - q_f$ and the same unit cost $c$. They
choose quantities at the same time (Cournot). Home's government pays its firm $s$
per unit sold; Home's welfare is the firm's profit minus the subsidy, $W = (P - c)q_h$.

\begin{proposition}[Export subsidy under quantity competition]\label{prop:b3:trade-policy:bs}
In the Cournot equilibrium $q_h = (a - c + 2s)/3$ and $q_f = (a - c - s)/3$. Home's
welfare is maximised by the subsidy $s^* = (a - c)/4$, which raises it from
$(a - c)^2/9$ to $(a - c)^2/8$; the Foreign firm's profit falls from $(a - c)^2/9$ to
$(a - c)^2/16$.
\end{proposition}

\begin{proof}
Home's firm maximises $(a - q_h - q_f - c + s)q_h$, so $q_h = (a - c + s - q_f)/2$; Foreign's
gives $q_f = (a - c - q_h)/2$. Solving, $q_h = (a - c + 2s)/3$, $q_f = (a - c - s)/3$, and
$P - c = (a - c - s)/3$. So $W(s) = (a - c - s)(a - c + 2s)/9$, and $W'(s) = [(a - c) - 4s]/9$
vanishes at $s^* = (a - c)/4$. Then $q_h = (a - c)/2$, $q_f = (a - c)/4$, $P - c = (a - c)/4$,
$W = (a - c)^2/8$ and Foreign's profit is $(a - c)^2/16$. At $s = 0$ both earn $(a - c)^2/9$.
\end{proof}

At $s^*$ the Home firm sells what it would sell if it could commit to its output
before its rival moved: the subsidy buys a commitment the firm cannot make alone.
The result turns on the form of competition. When firms set prices for
differentiated goods, prices are strategic complements: a Home firm that raises
its price leads its rival to raise its own. Home then gains from an export
\emph{tax}, which commits its firm to a higher price, not from a subsidy. The listing
computes both.

\omcode{code/bachelor-3/11-trade-policy/listings/strategic_trade.py}{1}{33}{Home's best export subsidy when the firms compete in quantities and when they compete in prices.}\label{lst:b3:trade-policy:strategic}

\omoutput{code/bachelor-3/11-trade-policy/outputs/strategic_trade.txt}{Welfare in units of the good's value.}\label{out:b3:trade-policy:strategic}

With $a = 10$ and $c = 2$, the best subsidy is 2 and Home's welfare rises from 7.11
to 8; under price competition the best policy is a tax of about 0.36, and Foreign's
profit rises with it. % output: strategic_trade
The sign of the right policy depends on a fact about conduct that a government
rarely knows. If both governments subsidise, they play the tariff game again, and
the rent goes to the buyers in the third market.

\begin{model}{Strategic export subsidies}
\assumes A few firms with market power share a world market; the government
moves first and can commit; the form of competition is known.
\predicts Under quantity competition, an export subsidy shifts profit to the home
firm and raises national welfare; under price competition, an export tax does.
\tested Not tested here: a test would need estimates of firms' conduct and of the
rent shifted in a subsidised industry, with the rival government's response.
\breaks When the form of competition is unknown, the sign of the policy is
unknown; when rivals subsidise too, the gain goes to foreign buyers; when entry
is free, the rent the subsidy shifts disappears.
\end{model}

\section{The political economy of protection}\label{sec:b3:trade-policy:political}

If tariffs usually lower national welfare, why are they levied? The \omterm{def:b3:factor-endowments:specific-factors}{specific-factors
model} of \cref{ch:b3:factor-endowments} gives one answer: the gains of protection
are concentrated on the owners and workers of the protected industry, who know
what is at stake and organise; the losses are spread thinly over every buyer. The
tariff is then a transfer obtained through the political process --- a form of rent
seeking, which \cref{ch:b3:interest-groups-and-regulation} studies with its own models.

The protection-for-sale model of Grossman and Helpman makes this precise. Each
organised industry offers the government contributions that depend on its tariff;
the government maximises the contributions plus $a$ times aggregate welfare, where
the weight $a$ measures how much it cares about voters at large. In equilibrium
the tariff on good $i$ satisfies
\begin{equation}\label{eq:b3:tp:pfs}
  \frac{t_i}{1 + t_i} = \frac{I_i - \alpha_L}{a + \alpha_L}\, \frac{z_i}{e_i},
\end{equation}
where $I_i$ is 1 if the industry is organised and 0 if not, $\alpha_L$ is the share
of the population that belongs to some lobby, $z_i$ is domestic output over imports
and $e_i$ is the elasticity of import demand. Organised industries get protection,
the more so when their output is large relative to imports and their import
demand inelastic, because there the tariff's deadweight loss is small; unorganised
ones get a small import subsidy.

\begin{model}{Protection for sale}
\assumes Organised industries bid for tariffs with contributions; the government
weighs contributions against aggregate welfare, with weight $a$; small country.
\predicts \Cref{eq:b3:tp:pfs}: tariffs rise with the ratio of output to imports in
organised industries and fall with it in unorganised ones.
\tested Goldberg and Maggi (1999) found the pattern of United States protection in
1983 consistent with the model's predictions, and estimated that the government
weighed welfare 50 to 88 times as much as contributions. % ledger: b3tp:gm.pfs
Such a large weight means the model, fitted, predicts tariffs close to zero, and
it does not explain why some industries organise and others do not.
\breaks When the organised industries are not given but formed by the stakes; when
tariffs are bound by agreements, as the next section shows.
\end{model}

\begin{record}[The tariff of 1930]
In 1930 the United States raised the duties of its tariff schedule, and as world
prices fell during the Depression, duties set in dollars per unit rose further
as a percentage of prices. The volume of United States imports fell by over 40~\%.
Irwin estimated that the tariff itself cut imports by 4--8~\%, and the interaction
of specific duties with deflation by another 8--10~\%: together, nearly a quarter
of the fall; the rest came from the collapse of income. % ledger: b3tp:irw.sh
Other countries raised their own barriers in the following years, the tariff game
of \cref{sec:b3:trade-policy:game} played for real.
\end{record}

\section{The rules of the trading system}\label{sec:b3:trade-policy:rules}

The tariff game has a cooperative outcome that neither player reaches alone.
Since 1947 most countries have bound themselves to rules, first in the General
Agreement on Tariffs and Trade and since 1995 in the World Trade Organization,
which turn the repeated game of \cref{exo:b3:trade-policy:7} into written
commitments.

\begin{definition}[Most-favoured-nation treatment; reciprocity; bound tariff]\label{def:b3:trade-policy:mfn}
\index{most-favoured-nation treatment}\index{reciprocity}\index{bound tariff}\index{binding overhang}
\emph{Most-favoured-nation} (MFN) treatment is the rule that a tariff applied to one
member's goods must be applied to every member's: no member may be favoured,
except within a free trade area or a customs union. \emph{Reciprocity} is the
practice of exchanging tariff cuts of similar value in negotiation. A \emph{bound
tariff} is the ceiling a member commits not to exceed; the \emph{binding overhang} is
the bound rate minus the applied rate.
\end{definition}

\begin{definition}[Anti-dumping duty; safeguard]\label{def:b3:trade-policy:antidumping}
\index{anti-dumping duty}\index{safeguard}
An \emph{anti-dumping duty} is a tariff on a foreign firm's good found to be sold
below its home price or its cost, and to injure the importing industry. A
\emph{safeguard} is a temporary tariff or quota on all imports of a good whose
surge injures the importing industry, without any finding of unfair pricing.
\end{definition}

The two rules have an economic reading. With MFN, a tariff cut offered to one
partner reaches all, so each member's tariff moves the world price of its imports
as one. \omterm{def:b3:trade-policy:mfn}{Reciprocity} then neutralises the \omterm{def:b3:trade-policy:optimal-tariff}{terms-of-trade effect}: when both sides cut
together, the world prices of their imports change little, and only the efficiency
gains remain. In this reading, a trade agreement exists to stop large countries
setting their \omterm{def:b3:trade-policy:optimal-tariff}{optimal tariffs}. Countries outside the agreement should then set
higher tariffs where their market power is larger.

\begin{record}[Optimal tariffs outside the agreement]
Broda, Limão and Weinstein estimated the elasticity of foreign export supply for
thousands of goods and countries that were not members of the World Trade
Organization. Those countries set tariffs about 9 percentage points higher on
goods supplied inelastically than on goods supplied elastically, a gap as large
as their average tariffs; and the United States' barriers not covered by the
agreement were higher where its market power was larger. % ledger: b3tp:blw.optimal
\end{record}

Bound rates are a ceiling, often far above the applied rates.

\begin{omchart}{Bound against applied tariff rates, simple means over all products,
each country's latest year from 2017 to 2021 (percent; the 27 members of the
European Union share one schedule and are one point). The line marks bound equal
to applied.}{World Bank, World Development Indicators (TM.TAX.MRCH.BR.ZS,
TM.TAX.MRCH.SM.AR.ZS); our computation}
\begin{tikzpicture}
\begin{axis}[omscatter, width=0.68\linewidth, height=0.48\linewidth, xmin=0, xmax=20,
  ymin=0, ymax=170, xlabel={Applied tariff (\%)}, ylabel={Bound tariff (\%)}]
  \addplot table[x=applied, y=bound] {figdata/bachelor-3/11-trade-policy/out/bound.csv};
  \addplot[only marks, mark=*, mark size=1.6pt, omThm, nodes near coords,
    point meta=explicit symbolic, every node near coord/.append style={font=\scriptsize, anchor=west}]
    table[x=applied, y=bound, meta=shown] {figdata/bachelor-3/11-trade-policy/out/bound_labels.csv};
  \addplot[omThm, thin, domain=0:20, samples=2] {x};
  \node[font=\scriptsize, omThm, anchor=west] (rich) at (axis cs:7,3) {EU, Japan, United States};
  \draw[omThm, thin] (rich.west) -- (axis cs:3.0,3.6);
\end{axis}
\end{tikzpicture}
\end{omchart}

% figdata: bound_stats, bound_labels
Across 131 schedules, the median bound rate is 29.7~\% and the median applied rate
7.7~\%; in 72~\% of them the overhang exceeds 10 points. Bangladesh has bound its
tariffs at 158.9~\% on average and applies 12.6~\%; the European Union, the United
States and Japan apply rates within about 2 points of their bindings. A large
overhang lets a government raise tariffs without breaking a commitment, so the
bindings of poorer members constrain them less than the opening chart suggests
(\cref{exo:b3:trade-policy:8}).

Anti-dumping is the rules' main escape route. Its cases have spread from a few
countries to many, with new users making up half of the world's total, and an
\omterm{def:b3:trade-policy:antidumping}{anti-dumping duty} cuts the value of the targeted imports by 30--50~\% on average.
% ledger: b3tp:prusa.ad
Disputes go to panels, and since 1995 members have filed 646 requests for
consultations, the first step of a dispute.
% ledger: b3tp:wto.ds

\begin{dated}{2026}
The Appellate Body, which hears appeals against panel reports, cannot review
appeals: the term of its last member expired on 30 November 2020, and its seats
are vacant. A panel report under appeal therefore stays without effect.
% ledger: b3tp:wto.ab
\end{dated}

\begin{debate}[Does membership raise trade?]
If the agreement removes the terms-of-trade temptation, members should trade more.
\begin{itemize}
\item Rose (2004) estimated a \omterm{def:b3:gravity:gravity}{gravity equation} on a panel of 175 countries over
fifty years and found little evidence that joining or belonging to the agreement
changed a country's trade; the system of tariff preferences for developing
countries was associated with about twice as much trade. % ledger: b3tp:rose.wto
\emph{Grade 3}: a large panel with gravity controls, but membership is not as good as
random, and the dummy treats very different commitments alike.
\item Subramanian and Wei (2007) let the effect differ between members: membership
raised the imports of industrial countries, which made deep cuts, by an amount
estimated at about 44~\% of world trade, and did little for developing members
that cut little, except those that joined after 1995. % ledger: b3tp:sw.wto
\emph{Grade 3}: the same design, with the effect allowed to follow the commitments.
\end{itemize}
Membership in name has no effect; the tariff cuts that membership brought do.
\end{debate}

\section{Evidence}\label{sec:b3:trade-policy:evidence}

\begin{definition}[Tariff pass-through]\label{def:b3:trade-policy:pass-through}
\index{tariff pass-through}
\emph{Tariff pass-through} is the share of a tariff that shows up in the price paid
by the importer, tariff included: complete pass-through means foreign exporters'
prices do not fall at all; zero means they absorb the whole tariff.
\end{definition}

\begin{record}[The U.S.--China tariffs of 2018--2019]
In 2018 the United States raised tariffs on a wide range of imports, many of them
from China, and its partners retaliated. Two studies of the import prices found
complete pass-through: the prices that foreign exporters received did not fall,
so the tariffs were paid by U.S. buyers. % ledger: b3tp:arw.passthrough, b3tp:fgkk.loss
One estimated the loss to U.S. buyers of imports at \$51 billion, 0.27~\% of GDP,
and the aggregate loss, after tariff revenue and gains to protected producers, at
\$7.2 billion, 0.04~\% of GDP. % ledger: b3tp:fgkk.loss
At the store, prices of affected goods rose less than at the border: retailers'
margins absorbed part of the tariff. And U.S. exporters hit by foreign retaliation
lowered their prices. % ledger: b3tp:cgnt.store
\end{record}

\begin{model}{The large-country tariff}
\assumes Competitive markets; an upward-sloping foreign export supply; no
retaliation.
\predicts Foreign exporters absorb part of a tariff through a lower world price,
and a small tariff raises the importing country's welfare; the \omterm{def:b3:trade-policy:optimal-tariff}{optimal tariff} is
$1/\varepsilon^*$.
\tested In 2018--2019, the world prices of the U.S. imports targeted by tariffs did not
fall: pass-through was complete, as if foreign export supply were perfectly
elastic, and no terms-of-trade gain appeared. U.S. exporters facing retaliation
did lower their prices. % ledger: b3tp:arw.passthrough, b3tp:cgnt.store
\breaks Retaliation, which turns the gain into a loss; exporters who can sell
elsewhere at the same price; imperfect competition.
\end{model}

\begin{debate}[Trade and wages]
How much did import competition lower wages and employment for the workers exposed
to it, and did the gains elsewhere offset it?
\begin{itemize}
\item Autor, Dorn and Hanson measured each U.S. local labour market's exposure to
Chinese imports from 1990 to 2007 by its initial mix of industries, instrumented
with Chinese exports to other rich countries. Exposed areas had higher
unemployment, lower participation and wages, and higher transfer payments; the
shock explains at least a quarter of the fall in manufacturing employment.
% ledger: b3tp:adh.china
\emph{Grade 2}: a strong shift-share design, whose validity rests on either the
industry shocks being as good as random or the initial shares being so --- the
two readings developed by Borusyak, Hull and Jaravel and by Goldsmith-Pinkham,
Sorkin and Swift. % ledger: b3tp:bhj.shiftshare, b3tp:gpss.bartik
\item Dauth, Findeisen and Suedekum applied the same design to Germany, 1988--2008:
regions specialised in import-competing industries lost jobs, but export-oriented
regions gained more, some 442{,}000 jobs in all, almost all from trade with
Eastern Europe rather than China. % ledger: b3tp:dfs.germany
\emph{Grade 2}, with the same caveats; the two studies disagree in sign for the
aggregate, because the two countries' export opportunities differed.
\item Caliendo, Dvorkin and Parro built a model of 38 countries, 50 states and 22
sectors: China's import competition cost about 1 million U.S. manufacturing jobs
and raised aggregate U.S. welfare, with large losses in some labour markets.
% ledger: b3tp:cdp.dynamics
\emph{Grade 4}: a calibrated model, whose aggregate gain rests on its assumptions
about mobility and trade.
\end{itemize}
\end{debate}

\begin{record}[Infant industries: the evidence]
The infant-industry argument holds that a young industry, protected for a time,
learns to produce at world costs and can then stand alone. The claim is hard to
test, because governments protect industries they expect to grow. Juhász found a
natural experiment: during the Napoleonic Wars (1803--1815), the blockade against
Britain cut some regions of the French Empire off from British trade more than
others, for reasons of geography. The better protected regions built more
mechanised cotton-spinning capacity, kept more of the industry in the long run,
and had higher industrial value added per head until the second half of the
nineteenth century, but not later. % ledger: b3tp:juh.blockade
One episode, in one industry with a new technology: \emph{grade 2} for that
setting, and silent on how to pick which industries to protect.
\end{record}

\section{Exercises}\label{sec:b3:trade-policy:exercises}

\begin{exercise}[$\star$]\label{exo:b3:trade-policy:1}
In the opening chart, which economy cut its average tariff most in points, and
which in percent? In which decade did China's fall most?
\end{exercise}

\begin{exercise}[$\star$]\label{exo:b3:trade-policy:2}
A small country's demand is $Q = 80 - P$ and its supply $Q = P - 20$; the world price
is 40. Find imports under free trade and with a specific tariff of 5, the tariff
revenue and the deadweight loss.
\end{exercise}

\begin{exercise}[$\star$]\label{exo:b3:trade-policy:3}
A car sells for 20{,}000 at world prices and uses inputs worth 12{,}000. Cars bear a
10~\% tariff and inputs 5~\%. Find the \omterm{def:b3:trade-policy:erp}{effective rate of protection}. What is it if
inputs bear 20~\%?
\end{exercise}

\begin{exercise}[$\star$]\label{exo:b3:trade-policy:4}
What price data would show that foreign exporters absorbed a tariff in full? Which
prices would you compare, and with what?
\end{exercise}

\begin{exercise}[$\star\star$]\label{exo:b3:trade-policy:5}
A large country's import demand is $M = 100 - 2P$ and foreign export supply $X^* = 2P^*
- 20$. It levies a specific tariff of 10. Find $P^*$, $P$ and $M$, the terms-of-trade
gain, the deadweight loss and the net welfare change.
\end{exercise}

\begin{exercise}[$\star\star$]\label{exo:b3:trade-policy:6}
Foreign export supply has a constant elasticity of 4. Find the \omterm{def:b3:trade-policy:optimal-tariff}{optimal tariff}.
What happens to it as the elasticity grows?
\end{exercise}

\begin{exercise}[$\star\star$]\label{exo:b3:trade-policy:7}
The tariff game of this chapter is played every year, with discount factor $\beta$.
Each country plays free trade as long as the other has, and tariffs for ever after
any tariff. For which $\beta$ is free trade an equilibrium?
\end{exercise}

\begin{exercise}[$\star\star$]\label{exo:b3:trade-policy:8}
Using the World Bank's bound and applied tariff series, compute each country's
\omterm{def:b3:trade-policy:mfn}{binding overhang} in its latest year from 2017 to 2021, counting the European Union
once. In what share of countries could the government double its applied tariff
without breaking its binding? Is the overhang larger in poorer countries?
\end{exercise}

\begin{exercise}[$\star\star$]\label{exo:b3:trade-policy:9}
A government protects a young industry and twenty years later it exports. Does that
prove that protection worked? What data would refute the claim?
\end{exercise}

\begin{exercise}[$\star\star\star$]\label{exo:b3:trade-policy:10}
Prove the optimal-tariff formula for a specific tariff $s$, with $P = P^* + s$: show
that at an interior optimum $s/P^* = 1/\varepsilon^*$.
\end{exercise}

\begin{exercise}[$\star\star\star$]\label{exo:b3:trade-policy:11}
Domestic demand is $P = 100 - Q$; the only domestic firm has marginal cost $20 + q$; the
world price is 40. Compare a specific tariff of 10 with the quota that gives the
same imports: domestic price, output and imports.
\end{exercise}

\begin{exercise}[$\star\star\star$]\label{exo:b3:trade-policy:12}
Five regions have the 1990 employment shares and national import shocks of the
chapter's exercise file. Compute each region's shift-share exposure and regress the
change in its manufacturing employment share on it. What must be true of the
shocks or the shares for the slope to be causal?
\end{exercise}

\section{Problem: Brenland's Steel Tariff}\label{sec:b3:trade-policy:problem}

\begin{problem}[{Weekend problem --- Brenland's steel tariff}]\label{pb:b3:trade-policy:steel}
In Brenland, steel demand is $100 - 2P$ and domestic supply $-20 + 2P$ (millions of
tonnes; $P$ in dalens per tonne). Kestria's export supply of steel is $X^* = -40 + 4P^*$.
Brenland exports grain to Kestria: Brenland's export supply of grain is $-20 + 2p$
and Kestria's import demand $100 - 3p_K$. Welfare is measured in millions of
dalens. Brenland considers a 25~\% tariff on steel.

\medskip\noindent\textbf{Part I --- the small-country benchmark.}
\begin{enumerate}
\item Find Brenland's import demand for steel.
\item Find the free-trade price and imports when Brenland and Kestria trade steel.
\item Suppose instead Brenland were small and the world price stayed at 20. Find
the domestic price, production, consumption and imports with the tariff.
\item Find the changes in consumer and producer surplus and the tariff revenue.
\item Find the deadweight loss and the net welfare change.
\item Why is a tariff always costly for a small country?
\end{enumerate}

\medskip\noindent\textbf{Part II --- the large-country tariff.}
\begin{enumerate}[resume]
\item With $P = 1.25P^*$, find the world price, the Brenland price and imports.
\item Who pays the tariff? Give the share of the price wedge borne by Kestrian
exporters.
\item Find the terms-of-trade gain and the efficiency loss.
\item Find Brenland's net welfare change, and check it against the sum of the
surplus changes and revenue.
\item Find Kestria's loss in the steel market.
\item Find the elasticity of Kestria's export supply at free trade and at the 25~\%
tariff. Is 25~\% above or below the \omterm{def:b3:trade-policy:optimal-tariff}{optimal tariff}?
\item Find the \omterm{def:b3:trade-policy:optimal-tariff}{optimal tariff} and Brenland's welfare gain at that tariff.
\end{enumerate}

\medskip\noindent\textbf{Part III --- retaliation.}
\begin{enumerate}[resume]
\item Find the free-trade grain price and Brenland's grain exports.
\item Kestria retaliates with a 25~\% tariff on Brenland's grain. Find the price
Brenland's farmers receive and the exports.
\item Find Brenland's loss in the grain market.
\item Find Kestria's gain in the grain market.
\item Write the tariff game in normal form, with welfare changes from free trade.
\item Find the Nash equilibrium. Is it efficient?
\item What does the 2018--2019 record say about the assumption of Part II?
\end{enumerate}

\medskip\noindent\textbf{Part IV --- steel users and jobs.}
\begin{enumerate}[resume]
\item Brenland's cars sell at world prices and are not protected; steel is 30~\% of
the value of a car at free-trade prices. Find the \omterm{def:b3:trade-policy:erp}{effective rate of protection} of
car making under the large-country tariff.
\item Steel employs 1{,}000 workers per million tonnes. How many steel jobs does
the tariff create?
\item Car making employs 50{,}000 workers. If employment moves in proportion to
value added, how many car jobs are lost?
\item Why do jobs counted this way not measure the tariff's welfare effect?
\item Which group in Brenland would lobby for the tariff, and which against it?
\item State the result: the net welfare change of the tariff for Brenland, with and
without retaliation.
\end{enumerate}
\end{problem}
